\mode*
\subsection{Dictionaries}
\label{dictionaries}

The literature has little to say about dictionaries. The one study we have
found of how students understand them is a think-aloud study of two students
learning Python \autocite{Stansbury2015Dictionaries}; the studies of hash
tables concern how a table is implemented, for instance how it is resized
\autocite{Karpierz2014HashTables}, not how a program uses one. The analyses
below therefore rest mainly on the misconceptions that the course material
itself targets, and each says what, if anything, the literature documents.

A dictionary is written with the same square brackets as a list:
\mintinline{python}{phone["ronja"]} looks like \mintinline{python}{names[0]}.
The likeness is useful, since the students can already predict what the
brackets do on a list, but it is also where the misconceptions come from.
The four analyses below share one source: a student who reads the dictionary
as a list --- with positions, growing only on request, failing quietly or not
at all --- predicts the dictionary wrongly in four different places. Each
analysis keeps the brackets invariant and varies what they hold or what is
done through them, so that the difference between the two containers, not
their shared notation, is what is discerned.
The analyses follow the course's lecture on dictionaries, whose programs
the code below abridges.

The first object of learning is \emph{how an entry of a dictionary is
reached}. The critical aspect is that it is reached by its \emph{key}, not by a
position: the entries of a dictionary lie in no order that can be counted with
0, 1, 2. The misconception it targets is that a dictionary is a list that
happens to accept strings as indices.
The closest documented case is \citeauthor{Stansbury2015Dictionaries}'s:
students who conflated Python's lists and dictionaries, one of whom
\textcquote{Stansbury2015Dictionaries}{treated her dictionary as if it were a
literal list of two-element lists and could be accessed as such}. Her list held
pairs rather than values indexed by strings, but in both readings an entry has
a position.

\begin{description}
  \item[Contrast] The same index in a list and in a dictionary holding the
    same two persons. The brackets and their contents are invariant; only the
    kind of container varies. The list gives its first element; the dictionary
    stops the program with \mintinline{python}{KeyError: 0}, since 0 is not a
    key in it --- there is no \emph{first} entry to give.

    \hfill
    \begin{minipage}[t]{0.45\columnwidth}
      \begin{minted}[highlightlines={2}]{python}
        names = ["ronja", "pippi"]
        print(names[0])
        # ronja
      \end{minted}
    \end{minipage}
    \hfill
    \begin{minipage}[t]{0.45\columnwidth}
      \begin{minted}[highlightlines={3}]{python}
        phone = {"ronja": "0701234567",
                 "pippi": "0721234567"}
        print(phone[0])
        # KeyError: 0
      \end{minted}
    \end{minipage}
    \newline
    The contrast is not complete until it is turned round: holding the list
    invariant and varying what stands in the brackets,
    \mintinline{python}{names["ronja"]} fails too, with a
    \mintinline{python}{TypeError}. That
    \mintinline{python}{phone["ronja"]} works is compatible with the
    misconception; that each container rejects the other's lookup is not.

  \item[Generalisation] Keeping the critical value invariant --- the brackets
    hold a key --- we vary the type of the keys. With integer keys the key
    \emph{looks} like an index, and still is not one: the lookup finds the
    entry whose key is 0, wherever it stands.

    \hfill
    \begin{minipage}[t]{0.55\columnwidth}
      \begin{minted}[highlightlines={2}]{python}
        floors = {1: "hall", 0: "cellar"}
        print(floors[0])
        # cellar
      \end{minted}
    \end{minipage}
    \hfill\null
\end{description}

The second object of learning is \emph{what an assignment to a key does}. The
critical aspect is that the same assignment does one of two things depending on
the key: to a key that is there it \emph{replaces} the value, to a key that is
not there it \emph{adds} an entry. The practical consequence is that a misspelt
key raises no error but quietly adds an entry of its own --- the fault
students then hunt for in their programs.
We have found no study documenting what students expect an assignment to a key
to do; the misconception is the one the course material itself targets.

\begin{description}
  \item[Contrast] We hold the assignment invariant and vary whether its key is
    already in the dictionary. The first assignment adds, the second replaces;
    nothing in the statement tells them apart, only the dictionary does.

    \hfill
    \begin{minipage}[t]{0.55\columnwidth}
      \begin{minted}[highlightlines={2,4}]{python}
        phone = {"ronja": "0701234567"}
        phone["pippi"] = "0721234567"
        # {'ronja': ..., 'pippi': ...}
        phone["ronja"] = "08-790 00 00"
        # {'ronja': '08-790 00 00', 'pippi': ...}
      \end{minted}
    \end{minipage}
    \hfill\null

  \item[Contrast] We then hold the assignment to a place that does not exist
    invariant and vary the container. The list does not grow by assignment:
    it stops the program with an \mintinline{python}{IndexError} and grows
    only when asked, with \mintinline{python}{append}. The dictionary grows by
    itself.

    \hfill
    \begin{minipage}[t]{0.45\columnwidth}
      \begin{minted}[highlightlines={2}]{python}
        phone = {"ronja": "0701234567"}
        phone["pippi"] = "0721234567"
        # adds an entry
      \end{minted}
    \end{minipage}
    \hfill
    \begin{minipage}[t]{0.45\columnwidth}
      \begin{minted}[highlightlines={2}]{python}
        numbers = ["0701234567"]
        numbers[1] = "0721234567"
        # IndexError
      \end{minted}
    \end{minipage}
    \newline
    Together the two contrasts explain the misspelt key:
    \mintinline{python}{phone["rnoja"] = ...} is an assignment to a key that is
    not there, so it adds.
\end{description}

The third object of learning is \emph{what happens when a key is missing}. The
critical aspect is that a lookup with the brackets treats a missing key as an
error, a \mintinline{python}{KeyError}, and that a program which expects
missing keys must say what it wants instead.
We have found no study of what students expect a lookup of a missing key to
give. That the error does stop novices' programs is documented: among nearly a
million novice Python programs on Python Tutor that crashed and were then
fixed, \mintinline{python}{KeyError} was one of the five most common uncaught
exceptions, though, at 3\,\%, the least common of them
\autocite{Cosman2020PABLO}. Those data show what the programs did, not what
their authors expected.
There is more than one way to say it, and the ways are values in one
dimension: how the program learns that the key is missing.

\begin{description}
  \item[Contrast] We hold the dictionary, the question to the user and the
    answer (a name that is not in the dictionary) invariant, and vary how the
    missing key is handled: test first with \mintinline{python}{in}, look up
    and catch the \mintinline{python}{KeyError}, or ask
    \mintinline{python}{get} for a fallback. The three values are shown
    side by side rather than one at a time, since it is their simultaneous
    difference against the invariant program that lets the dimension itself
    be discerned \autocite[pp.~45--47]{NCOL}.

    \hfill
    \begin{minipage}[t]{0.3\columnwidth}
      \begin{minted}[fontsize=\footnotesize]{python}
        if name in phone:
            print(phone[name])
        else:
            print("unknown")
      \end{minted}
    \end{minipage}
    \hfill
    \begin{minipage}[t]{0.3\columnwidth}
      \begin{minted}[fontsize=\footnotesize]{python}
        try:
            print(phone[name])
        except KeyError:
            print("unknown")
      \end{minted}
    \end{minipage}
    \hfill
    \begin{minipage}[t]{0.3\columnwidth}
      \begin{minted}[fontsize=\footnotesize]{python}
        print(phone.get(
            name, "unknown"))
      \end{minted}
    \end{minipage}
    \newline
    All three print the same; they differ in what they say about the program
    --- whether a missing key is an expected case or something that should not
    happen. Without a fallback, \mintinline{python}{get} returns
    \mintinline{python}{None}, the value the brackets never give: they give
    the value or an error.
\end{description}

The fourth object of learning is \emph{what a loop over a dictionary goes
through}. The critical aspect is that the loop variable gets the \emph{keys},
and nothing else; the values and the key--value pairs must be asked for.
We have found no study of what students expect such a loop to give. The
list-of-pairs reading that \textcite{Stansbury2015Dictionaries} observed would
predict the pairs, but that study did not ask about loops.

\begin{description}
  \item[Contrast] We hold the dictionary and the form of the loop invariant and
    vary only what the loop goes through: the dictionary itself, its
    \mintinline{python}{values()}, or its \mintinline{python}{items()}, which
    gives a tuple at a time, unpacked in the loop header as with
    \mintinline{python}{enumerate}.

    \hfill
    \begin{minipage}[t]{0.3\columnwidth}
      \begin{minted}{python}
        for x in phone:
            print(x)
        # ronja
        # pippi
      \end{minted}
    \end{minipage}
    \hfill
    \begin{minipage}[t]{0.3\columnwidth}
      \begin{minted}{python}
        for x in phone.values():
            print(x)
        # 0701234567
        # 0721234567
      \end{minted}
    \end{minipage}
    \hfill
    \begin{minipage}[t]{0.3\columnwidth}
      \begin{minted}{python}
        for k, v in phone.items():
            print(k, v)
        # ronja 0701234567
        # pippi 0721234567
      \end{minted}
    \end{minipage}
\end{description}

The four analyses come together when a dictionary is built up rather than
written out, as when counting words. Predicting the result requires attending
at once to the key as the way in (the word, not a position), to assignment
adding a new key and replacing an old one's count, to the missing key that
\mintinline{python}{get} turns into a zero instead of an error, and to the loop
that hands out the words.

\begin{description}
  \item[Fusion] The same line counts the first occurrence of a word and every
    later one: for a new word \mintinline{python}{get} gives the fallback and
    the assignment adds the key, for a known word it gives the count and the
    assignment replaces it. The loop then goes through the keys.

    \hfill
    \begin{minipage}[t]{0.8\columnwidth}
      \begin{minted}[highlightlines={3}]{python}
        counts = {}
        for word in ["to", "be", "or", "not", "to", "be"]:
            counts[word] = counts.get(word, 0) + 1
        for word in counts:
            print(word, counts[word])
        # to 2
        # be 2
        # or 1
        # not 1
      \end{minted}
    \end{minipage}
    \hfill\null
\end{description}
